% Final section; CSV-derived claims and published metadata were audited
% separately so their evidence bases remain explicit.

\section{Dataset Description}
\label{sec:c}

\subsection{Data source}

The group uses the \textbf{``Predict Students' Dropout and Academic
Success''} dataset from the UCI Machine Learning Repository (ID~=~697),
published by Realinho, Vieira Martins, Machado, and Baptista
\citep{realinho2021dataset}, licensed CC~BY~4.0. The dataset is described
in detail in the accompanying paper \citep{realinho2022predicting}.

The data was collected at \textbf{Instituto Polit\'ecnico de Portalegre},
Portugal, spanning \textbf{17 undergraduate degree programs} (Agronomy,
Design, Nursing, Journalism and Communication, Management, Social
Service, Biofuel Production Technologies, Informatics Engineering, etc.),
academic years \textbf{2008/09 to 2018/19}. It combines information from
several administrative sources at enrollment time (admission records,
socio-economic data) with academic outcomes at the end of each of the
first two semesters \citep[Section~2, ``Data Description'']{realinho2022predicting}.

\subsection{Dataset description --- numbers confirmed from actual data}

The group did not copy figures from the UCI metadata page; instead it
directly inspected the raw dataset to confirm:

\begin{itemize}
  \item The raw file contains \textbf{4,424 rows and 37 columns}: 4{,}424
    samples (students), 36 features plus 1 label column \texttt{Target}.
  \item There are \textbf{zero} missing values across the whole dataset
    (the authors pre-cleaned the data).
  \item \textbf{Notable finding:} the UCI page lists 36 features, but
    Table~1 of the original paper lists only 34. Cross-checking the
    actual column list against that table, the discrepancy is exactly
    two columns: \texttt{Previous qualification (grade)} and
    \texttt{Admission grade}, both present in the released CSV but absent
    from the paper's table. That discrepancy is only visible by reading
    the actual column list, not by copying the metadata page.
\end{itemize}

The 36 features fall into six groups by nature: demographics (6 columns,
e.g.\ Gender, Age at enrollment), socio-economic status (8 columns, e.g.\
Tuition fees up to date, Scholarship holder), macroeconomics (3 columns:
Unemployment rate, Inflation rate, GDP), academic background at
enrollment (7 columns, e.g.\ Course, Admission grade), first-semester
outcomes (6 columns), and second-semester outcomes (6 columns).

\textbf{Target variable} has 3 classes, moderately imbalanced:

\begin{table}[H]
  \centering
  \caption{Target class distribution (n = 4{,}424)}
  \begin{tabular}{@{}lrrp{6cm}@{}}
    \toprule
    Class & Count & \% & Meaning \\
    \midrule
    Graduate & 2{,}209 & 49.9\% & Graduated on the standard timeline (3 years; 4 for Nursing) \\
    Dropout  & 1{,}421 & 32.1\% & Dropped out \\
    Enrolled & 794     & 17.9\% & Still enrolled past the standard timeline \\
    \bottomrule
  \end{tabular}
  \\[0.3em]
  \footnotesize Percentages are rounded to one decimal place and therefore
  sum to 99.9\% rather than 100\%; the unrounded values (49.93\%, 32.12\%,
  17.95\%) sum to 100\%.
\end{table}

\textbf{Enrolled} is both the minority class and the hardest to classify:
across every model configuration compared in Section~\ref{sec:g}, its
recall is the lowest of the three classes (e.g.\ 0.3836 for the baseline
M0), which is why the class-imbalance improvement (Section~\ref{sec:f},
f.2) targets this class.

\begin{figure}[H]
  \centering
  \includegraphics[width=0.7\textwidth]{A_target_distribution.png}
  \caption{Target class distribution (n = 4{,}424).}
  \label{fig:c1}
\end{figure}

Several of the observations below refer to how strongly a feature drives
the baseline decision tree's actual splits, measured by Gini importance
(the total reduction in class impurity attributed to that feature across
the whole tree, described further in Section~\ref{sec:d}). Table~\ref{tab:c-importance}
lists the top 10 features by this measure, plus one lower-ranked feature
included for contrast with the discussion below.

\begin{table}[H]
  \centering
  \caption{Top-10 features by Gini importance in the baseline tree (M0), plus one feature shown for contrast}
  \label{tab:c-importance}
  \begin{tabular}{@{}rlr@{}}
    \toprule
    Rank & Feature & Gini importance \\
    \midrule
    1  & Curricular units 2nd sem (approved)    & 0.340 \\
    2  & Curricular units 2nd sem (grade)       & 0.045 \\
    3  & Tuition fees up to date                & 0.045 \\
    4  & Curricular units 2nd sem (enrolled)    & 0.045 \\
    5  & Admission grade                        & 0.044 \\
    6  & Course                                 & 0.044 \\
    7  & Curricular units 1st sem (evaluations) & 0.034 \\
    8  & Previous qualification (grade)         & 0.031 \\
    9  & Curricular units 1st sem (approved)    & 0.029 \\
    10 & Mother's occupation                    & 0.028 \\
    \multicolumn{3}{c}{\ldots} \\
    24 & Scholarship holder                     & 0.010 \\
    \bottomrule
  \end{tabular}
\end{table}

Descriptive statistics from exploratory data analysis illustrate how
strongly some features relate to the target: among
students who had \textbf{not paid tuition on time}, 86\% ended up in the
Dropout class, versus 25\% Dropout / 56\% Graduate among students who had
paid on time. Students with a \textbf{scholarship} have a Dropout rate
26.5 percentage points lower than those without (12.2\% vs.\ 38.7\%);
their rate is about one-third as high.

\begin{figure}[H]
  \centering
  \includegraphics[width=0.9\textwidth]{A_tuition_scholarship_vs_target.png}
  \caption{Target distribution split by \texttt{Tuition fees up to date}
    (left) and \texttt{Scholarship holder} (right).}
  \label{fig:c2}
\end{figure}

\texttt{Tuition fees up to date} is among the top-5 features by
permutation importance in the original paper \citep{realinho2022predicting},
and is also rank~3 in Table~\ref{tab:c-importance}: the two ways of
measuring importance agree here. \texttt{Scholarship holder} is a more
subtle case: it is \textbf{not} in the top-5 of either the paper or
Table~\ref{tab:c-importance} (rank~24), even though its
percentage-of-target chart above looks just as striking as tuition's. A
large gap in a simple percentage-of-target chart does not automatically
mean a feature drives an important tree split once it sits alongside
stronger, possibly overlapping features such as \texttt{Tuition fees up
to date}; the tree may instead be extracting the same signal from that
stronger correlate.

By program of study, the two lowest-graduation programs are course code
33 (Biofuel Production Technologies, $n=12$: 8 Dropout, 3 Enrolled, and
1 Graduate; 66.7\% Dropout) and course
code 9119 (Informatics Engineering, $n=170$, 54.1\% Dropout). Only
\textbf{8.3\% and 8.2\%} of their students graduate respectively, versus
the dataset average of 49.9\%. Course 33's rate is based on a very small
sample ($n=12$) and is not evidence of a stable program-level effect. The
two highest-graduation programs are course code 9238 (Social Service,
$n=355$, 69.9\% Graduate) and course code 9500 (Nursing, $n=766$, 71.5\%
Graduate), with the lowest Dropout rates (18.3\% and 15.4\%).

\begin{figure}[H]
  \centering
  \includegraphics[width=0.85\textwidth]{A_target_by_course.png}
  \caption{17 programs sorted by increasing Graduate rate.}
  \label{fig:c3}
\end{figure}

The original paper ranks \texttt{Course} in its top-5 features by
permutation importance \citep{realinho2022predicting}; in
Table~\ref{tab:c-importance}, \texttt{Course} ranks 6th, close to top-5
either way. Under the study's feature-availability assumption, it is treated
as available at enrollment and therefore retained in M3
(Section~\ref{sec:f}, f.3).

The strongest single predictor in the whole dataset, by contrast, is
only known once a semester has actually happened:

\begin{figure}[H]
  \centering
  \includegraphics[width=0.7\textwidth]{A_curricular_units_2nd_sem_by_target.png}
  \caption{\texttt{Curricular units 2nd sem (approved)} by Target class.}
  \label{fig:c4}
\end{figure}

Median values are 0 for Dropout, 4 for Enrolled, and 6 for Graduate. This
is the single most important feature in Table~\ref{tab:c-importance}
(rank~1, importance 0.340, more than the sum of ranks 2--5 combined), but
it is defined as a second-semester outcome and is therefore excluded from M3.

\subsection{Preprocessing performed}

All preprocessing logic is implemented in a single shared pipeline, used
by the whole team so every model is evaluated on the exact same
train/test split.

\begin{enumerate}
  \item \textbf{No missing-value handling needed}: the dataset is
    already clean at the source (0 missing values).
  \item \textbf{Re-encoding categorical columns stored as integers.} Many
    columns are integer codes for unordered categories (e.g.\
    \texttt{Course} code 33 = Biofuel Production Technologies, code 9500
    = Nursing \citep[Variables Table, row ``Course'']{realinho2021dataset};
    these codes carry no ordinal relationship). Left as-is,
    scikit-learn would treat them as continuous and produce meaningless
    splits like \texttt{Course} $\le$ 8.5. The group used a mixed
    strategy:
  \begin{itemize}
    \item \textbf{One-hot encode} the 4 low-cardinality columns:
      \texttt{Marital Status} (6), \texttt{Application mode} (18),
      \texttt{Course} (17), \texttt{Previous qualification} (17), so
      each split reads as a clear yes/no question (e.g.\ ``is this the
      Nursing program?'').
    \item \textbf{Keep the raw integer code} for high-cardinality
      columns: the mother's and father's occupations (32/46 values),
      the mother's and father's qualifications (29/34 values), and
      \texttt{Nacionality} (21 values). This was a pragmatic choice to
      limit the number of sparse indicators. It imposes an artificial
      order on nominal codes, so thresholds on these columns are an
      acknowledged representation limitation and require cautious
      interpretation.
    \item After this step the feature count grows from 36 to
      \textbf{90} columns.
  \end{itemize}
  \item \textbf{No feature scaling.} A decision tree splits on a
    threshold of one feature at a time, not on distances between data
    points as KNN or linear regression do, so it is invariant to any
    monotonic per-column transform such as scaling. Adding
    \texttt{StandardScaler}/\texttt{MinMaxScaler} here would be
    redundant and would not change the learned tree structure at all.
  \item \textbf{Stratified train/test split}: 80\%/20\%, stratified by
    class to preserve the $\approx$50/32/18 class ratio in both sets.
    This guarantees that both partitions retain the observed class
    proportions. A fixed random seed is used so all five models in the
    project are evaluated on the exact same split.
\end{enumerate}

After preprocessing: \textbf{3{,}539 training samples / 885 test
samples}, each with 90 numeric features, ready to feed directly into a
decision tree classifier.

\subsection{Multicollinearity among features}

\begin{figure}[H]
  \centering
  \includegraphics[width=\textwidth]{A_correlation_heatmap.png}
  \caption{Pearson correlation heatmap across all numeric features.}
  \label{fig:c5}
\end{figure}

The heatmap's colors show the overall pattern, but individual cell
values are not printed on a 36$\times$36 grid; Table~\ref{tab:c-corr}
lists the exact values for the pairs discussed here so they can be
checked directly.

\begin{table}[H]
  \centering
  \caption{Highest Pearson correlations in the dataset}
  \label{tab:c-corr}
  \begin{tabular}{@{}llr@{}}
    \toprule
    Feature & Correlated with & $r$ \\
    \midrule
    Curricular units 1st sem (credited) & Curricular units 2nd sem (credited) & 0.945 \\
    Curricular units 1st sem (enrolled) & Curricular units 2nd sem (enrolled) & 0.943 \\
    Mother's occupation                 & Father's occupation                 & 0.910 \\
    Curricular units 1st sem (approved) & Curricular units 2nd sem (approved) & 0.904 \\
    Nacionality                         & International                       & 0.791 \\
    \bottomrule
  \end{tabular}
\end{table}

The three highest-correlation pairs in the whole dataset are
\texttt{1st sem credited} $\leftrightarrow$ \texttt{2nd sem credited},
\texttt{1st sem enrolled} $\leftrightarrow$ \texttt{2nd sem enrolled},
and \textbf{\texttt{Mother's occupation} $\leftrightarrow$
\texttt{Father's occupation}}, which is \emph{higher} than the
semester-approved pair (\texttt{1st/2nd sem approved}). The pair
\texttt{Nacionality} $\leftrightarrow$ \texttt{International} is still a
high correlation (international students almost always have a
non-Portuguese \texttt{Nacionality} code), but lower than the four
pairs above.

A decision tree is not as sensitive to multicollinearity as linear
regression (it does not need redundant features removed to converge),
but multicollinearity still splits feature importance across correlated
columns and can make the tree structure less stable across samples. These
correlations therefore provide diagnostic context, but they are not the
selection rule for M3. The canonical early-warning experiment removes
exactly the 12 semester-outcome columns under the feature-availability
assumption stated in Section~\ref{sec:f3}, while retaining both \texttt{Nacionality} and
\texttt{International} so that feature timing is the only intentional
difference from M0.

\subsection{Why this dataset is appropriate for decision tree modeling}

\begin{enumerate}
  \item \textbf{Features have clear, real-world semantics} (tuition,
    scholarship, semester outcomes, program of study, \ldots), so every
    split translates into an ordinary sentence. This plays directly to a
    decision tree's interpretability strength, in a way that image or
    text features cannot match as easily.
  \item \textbf{A mix of categorical and numerical features}: after
    nominal categories are encoded, the tree handles the resulting
    numerical and indicator columns without scaling, unlike KNN and many
    linear-model workflows, where feature scale directly affects
    distances or optimization.
  \item \textbf{A genuine class imbalance} (50\%/32\%/18\%): creates a
    realistic scenario for applying and comparing imbalance-handling
    techniques, one of the improvement directions the brief itself
    suggests (Section 3.2.d, ``Handling class imbalance'').
  \item \textbf{Documented feature families support a timing experiment}:
    the published documentation distinguishes semester outcomes from
    enrollment-related data. The project therefore evaluates an early-warning
    exclusion rule, while treating the exact availability of every retained
    field as an assumption to be verified from a data dictionary and snapshot
    timestamps before deployment.
  \item \textbf{No missing values}: lets the group spend its time on
    analysis and interpretation instead of data cleaning, matching the
    brief's own emphasis on clearly explaining how the tree is built and
    why the improvements help, beyond simply running a ready-made model.
\end{enumerate}

\subsection{Ethical note}

The dataset contains sensitive variables: gender, nationality, marital
status, and parents' education/occupation. Any tree splits on these
variables only reflect \textbf{correlations observed in one Portuguese
institution's sample from academic years 2008/09--2018/19
\citep[Section~2]{realinho2022predicting}}, not causal relationships, and
should not be used as grounds for any discriminatory policy toward
students in practice.
